\subsection{Training Signal across Evolution}
\label{sec:analysissignal}

We audit the saved RL reward groups from the original 4B experiments and
\texttt{Codex-Vanilla}. A task provides within-group reward signal when its
eight valid rollouts receive nonconstant rewards; we report the fraction of
such tasks among complete, valid groups in each round
(Figure~\ref{fig:researchdiagnostics}b). This measures the relative reward
contrast used by GRPO, retaining partial-credit rewards and excluding
incomplete groups rather than treating them as failures.

\textbf{Training success can exhaust the learning signal.}\quad
On $\tau^2$, every \texttt{Codex-Vanilla} task in R9 and R10 receives eight
rewards of one, leaving no within-group reward contrast. On BFCL, only
2 of 52 tasks retain reward variation in each of these rounds.
In comparison, \method{} retains signal on 70.5\% of valid $\tau^2$ tasks
and 35.5\% of BFCL tasks in its final recorded round, R9.
These observations motivate linking generated tasks to their actual training
outcomes so the researcher can detect reward saturation and refresh its
targets. They do not isolate the causal effect of context refinement.

\subsection{Environment Implementation Quality}
\label{sec:analysisenvironment}

\todo{Audit 20 sampled environments per benchmark, with eight rollouts each,
using an independent LLM judge to flag implementation, interaction, and verifier
defects. Replace the illustrative values in Figure~\ref{fig:researchdiagnostics}c
with measured environment-level audit pass rates (no defects flagged in the
inspected rollouts).}
